\section*{APPENDIX}
\setcounter{subsection}{0}
\renewcommand{\thesubsection}{\Alph{subsection}}
\label{app:proofs}

\subsection{Optimal Window Error}
\label{app:window-laws}

\emph{The best prediction for a fixed window.} Let $\mathcal R$ be the
closure of states reachable before a window $s=(a_1,\ldots,a_d)$, and
write $\Phi_s=\Phi_{a_d}\cdots\Phi_{a_1}$. Under
Section~\ref{sec:criterion}'s assumptions, $\mathcal R$ is compact and
independent of $s$. For scalar output, define
\[
 y_-(s)=\min_{\rho\in\mathcal R}\mathcal C\Phi_s(\rho),\quad
 y_+(s)=\max_{\rho\in\mathcal R}\mathcal C\Phi_s(\rho).
\]
Any prediction must be at least half this range away from one endpoint.
The midpoint $m_d(s)=[y_-(s)+y_+(s)]/2$ achieves this bound. Taking the
worst case over windows gives the optimal error $D_d/2$. Using the
reachable closure does not change this supremum.

\emph{Why QELMs can approach the optimum.} Continuity of the channels
on the compact window--state domain makes $y_-$, $y_+$, and $m_d$
continuous. Uniform QELM density therefore supplies a predictor within
any $\varepsilon>0$ of $m_d$, with total error at most
$D_d/2+\varepsilon$. This proves an infimum; a finite model need not
attain it exactly.

For vector outputs, the triangle inequality still gives the lower bound
$D_d/2$. Replaying the window from any fixed reachable state gives a
continuous predictor with error at most $D_d$. Approximating this replay
function proves that arbitrarily accurate emulation is possible exactly
when $D_d\to0$.

\emph{Randomness cannot recover missing history.} Two histories with
the same window induce the same prediction distribution. For their
scalar outputs $y_0,y_1$,
\[
 |y_1-y_0|\leq\mathbb E|y_0-\widehat y|+\mathbb E|y_1-\widehat y|.
\]
At least one expected absolute error is therefore at least half the gap.
Randomized predictions cannot lower the worst-case information floor.

\subsection{Observable Memory and Required Context}
\label{app:observable-proof}

\emph{Keep only differences that can affect an output.} Let $\Delta$
be the compact set of reachable state differences, including limits,
and let $V=\operatorname{span}_{\mathbb R}\Delta$. A difference is
irrelevant if no future input sequence can reveal it. Remove these
differences by defining
\[
 \mathcal N=\{v\in V:\mathcal C\Phi_w(v)=0\text{ for every finite word }w\}.
\]
The empty word is included. On $W=V/\mathcal N$, the induced update
and output maps are $A_a[v]=[\Phi_a(v)]$ and
$\overline{\mathcal C}[v]=\mathcal C(v)$. This is the
reachable--observable reduction used in
Section~\ref{sec:observable-forgetting}, specializing
\cite[Theorem~III.1]{petreczky2026stability}.

\emph{Measure memory by its possible future effect.} On $W$, define
\[
 \|x\|_{\rm obs}=\sup_w\|\overline{\mathcal C}A_wx\|.
\]
The invariant reachable set makes this finite, and removing $\mathcal N$
makes it a norm. Let $c_d=\sup_{|s|=d}\|A_s\|_{\rm obs}$, using the
induced operator norm.

Appending common inputs preserves reachability, so
$\Phi_w(\Delta)\subseteq\Delta$ for every word $w$.
Absorbing the first $\ell$ inputs of a length-$(d+\ell)$ word into
its starting difference therefore gives $D_{d+\ell}\leq D_d$.
For $v\in\Delta$, $|s|=d$, and any future word $w$ of length $\ell$,
\[
 \|\overline{\mathcal C}A_wA_s[v]\|\leq D_{d+\ell}\leq D_d.
\]
Taking suprema gives the upper bound $D_d$. The empty future word
recovers $\sup_{v\in\Delta,\,|s|=d}\|\overline{\mathcal C}A_s[v]\|=D_d$,
establishing equality.

The convex hull of the reachable differences
in $W$ contains a norm ball and is bounded. Thus constants
$\kappa_-,\kappa_+>0$, independent of $d$, satisfy
\[
 \kappa_-c_d\leq D_d\leq\kappa_+c_d.
\]
This comparison connects observable output gaps to the size of products
of the reduced updates.

\emph{Decay rate and window length.} The product norms satisfy
$c_{d+\ell}\leq c_dc_\ell$. Hence $D_d\to0$ implies $c_{d_0}<1$ for
some $d_0$, which forces exponential decay. Conversely, $r_{\rm vis}<1$
gives exponential decay of $c_d$ and therefore of $D_d$. Taking $d$th
roots in the comparison above yields
\[
 \lim_{d\to\infty}D_d^{1/d}=r_{\rm vis}.
\]
For $0<r_{\rm vis}<1$, bounding $D_d$ between $r_-^d$ and $r_+^d$ for
sufficiently large $d$, where $0<r_-<r_{\rm vis}<r_+<1$, and solving
$D_d<2\varepsilon$ gives
\[
 \lim_{\varepsilon\downarrow0}\frac{d_\varepsilon}{\log(1/\varepsilon)}
 =\frac{1}{-\log r_{\rm vis}}.
\]
If $W=\{0\}$, set $D_d=r_{\rm vis}=0$. The argument assumes freely
concatenable inputs; constrained input grammars require a separate
analysis.

\subsection{An Explicit QELM Construction}
\label{app:bernstein-resources}

\emph{Assumptions.} For a quantitative bound, take inputs in $[0,1]$,
scalar output $y=\operatorname{Tr}(O\rho)$, and $M=\|O\|_\infty$.
In addition to the assumptions above, suppose
\begin{equation}
 D_{\rm tr}(\Phi_a(\rho),\Phi_{a'}(\sigma))
 \leq\eta D_{\rm tr}(\rho,\sigma)+L_{\rm in}|a-a'|,
 \label{eq:incremental-lipschitz}
\end{equation}
with $0\leq\eta<1$. This is stronger than observable forgetting: it
requires uniform state contraction and Lipschitz dependence on the input.

\emph{Encode, measure, and average.} Encode each input $a_i$ in $r$
copies of $\sqrt{1-a_i}|0\rangle+\sqrt{a_i}|1\rangle$. Measuring each
block gives independent counts $K_i\sim\operatorname{Binomial}(r,a_i)$.
Assign the linear-readout coefficient $m_d(k_1/r,\ldots,k_d/r)$ to
count outcome $(k_1,\ldots,k_d)$. The expected prediction is the
Bernstein approximation
\[
 f_r(s)=\mathbb E\bigl[m_d(K_1/r,\ldots,K_d/r)\bigr].
\]
The midpoint inherits the input sensitivity bound
\[
 |m_d(s)-m_d(s')|
 \leq2ML_{\rm in}\sum_{i=1}^d\eta^{d-i}|a_i-a_i'|.
\]
Together with $\mathbb E|K_i/r-a_i|\leq1/(2\sqrt r)$, this bounds
the representation error. Averaging $Q$ independent readouts gives
$\widehat f_r$. Since each readout lies in $[-M,M]$, Hoeffding's
inequality~\cite{hoeffding1963probability} bounds its sampling error.
For each fixed history and $0<\zeta<1$, with probability at least
$1-\zeta$,
\begin{equation}
 |y(h)-\widehat f_r(S_d(h))|
 \leq D_d/2+\varepsilon_{\rm rep}+\varepsilon_{\rm shot},
 \label{eq:midpoint-resource-bound}
\end{equation}
where
\[
 \varepsilon_{\rm rep}=\frac{ML_{\rm in}(1-\eta^d)}{(1-\eta)\sqrt r},\quad
 \varepsilon_{\rm shot}=M\sqrt{\frac{2\log(2/\zeta)}{Q}}.
\]
With exact expectations, omit $\varepsilon_{\rm shot}$ and the bound
holds uniformly over histories. Contraction also gives
$D_d/2\leq M\eta^d$. The finite-shot guarantee is pointwise, not
simultaneous over all histories.

The construction uses $rd$ qubits and up to $(r+1)^d$ midpoint
coefficients. Computing these coefficients requires reachable-output
extrema and is not included in the resource count. This establishes
representability, not an efficient training procedure.

\emph{Analytical checks.} For the one-qubit channel in
Section~\ref{sec:sharpness}, the compatible populations have width
$D_d=\eta^d$ and midpoint
\[
 m_d(s)=\frac{\eta^d}{2}+(1-\eta)\sum_{i=1}^d\eta^{d-i}g(a_i).
\]
Thus any predictor $f$ has worst-case error
$\eta^d/2+\|f-m_d\|_\infty$. For $g(a)=a^2$, replacing $g(K_i/r)$ by
$K_i(K_i-1)/[r(r-1)]$ gives an unbiased readout for every $r\geq2$,
attaining the information floor exactly. For
$g(a)=(e^{2a}-1)/(e^2-1)$, the midpoint coefficients give a nonnegative
Bernstein bias. Additivity reduces the worst-case calculation to a
one-dimensional maximization of that bias. These checks use analytical
readouts and binomial expectations, not training.

\subsection{Protected Memory}
\label{app:protected-calibration}

\subsubsection{The Persistent Error Floor}
\label{app:protected-pulse}

Starting from $|00\rangle$, a first-qubit pulse writes singlet population
\[
 \delta=|\langle\psi^-|R_x^{(1)}(\theta)|00\rangle|^2
 =\tfrac12\sin^2(\theta/2).
\]
Without the pulse, this population is zero. Symmetric evolution with
collective loss preserves both populations during any unpulsed
continuation. The output gap therefore remains $\delta$ after arbitrarily
many shared inputs, giving $D_d\geq\delta$ and worst-case window error
at least $\delta/2$ for every finite $d$. A one-radian pulse gives
$\delta\approx0.1149$ and an error floor of approximately $0.0575$.
Conservation is needed only along these common continuations, not for
every possible input sequence.

\subsubsection{Dynamics and Controls}
\label{app:protected-implementation}

The numerical check uses
\[
 \begin{gathered}
 H(u)=JX_1X_2+0.37(Z_1+Z_2)+0.64u(X_1+X_2),\\
 L_c=\sqrt{0.8}(\sigma_1^-+\sigma_2^-),\qquad\Delta t=0.2.
 \end{gathered}
\]
Here $\sigma_i^-=|0\rangle\langle1|_i$ is the lowering operator.
Twelve independent realizations sample $J\in[-1,1]$ and 80-input
continuations with $u\in[0,1]$ after the optional pulse. Checks verify
the pulse value and singlet conservation. Matched local amplitude
damping reduces the median 80-step gap to $1.57\times10^{-5}$.
Relative collective-jump asymmetries of $0.02$, $0.05$, $0.10$, and
$0.20$ break exact conservation but need not erase the gap rapidly.
These controls distinguish protected memory from slow forgetting.
